% !TeX spellcheck = ru_RU
% !TEX root = vkr.tex

\section{Эксперимент}
\label{sec:experiment}

\subsection{Условия и исследовательские вопросы}

Проверка результата воспроизводилась \DTMdate{2026-05-27} на компьютере с
macOS (Darwin~25.4.0), AppleClang~21.0.0 и \CMake{}~4.3.2. Проект
компилировался стандартной конфигурацией \CMake{} с флагами
\texttt{-O2}, \texttt{-fsanitize=address,undefined} и
\texttt{-fno-omit-frame-pointer}, то есть с включёнными \AddressSanitizer{} и
\UndefinedBehaviorSanitizer{}; в непрерывной интеграции дополнительно
предусмотрена проверка тестов под \Valgrind{}. Для проверки были сформулированы
следующие вопросы:
\begin{description}
    \item[RQ1:] восстанавливает ли утилита исходные данные без потерь для
          типовых и граничных входов?
    \item[RQ2:] как распределение байтов во входном файле влияет на
          коэффициент сжатия статическим кодированием Хаффмана?
    \item[RQ3:] какое время требуется на сжатие и распаковку файлов выбранного
          размера в подготовленном сценарии измерений?
\end{description}

\subsection{Проверка корректности}

Модульные тесты проверяют построение дерева, чтение и запись заголовка,
побитовый ввод-вывод, основной кодек, случайные данные и обработку повреждённого
входа. Интеграционный сценарий последовательно сжимает и распаковывает
тестовые файлы, а затем сравнивает восстановленные данные с исходными. Итоги
запуска приведены в таблице~\ref{tab:tests}.

\begin{table}[htbp]
    \centering
    \caption{Результаты проверки корректности реализации.}
    \label{tab:tests}
    \begin{tabularx}{\textwidth}{@{}>{\raggedright\arraybackslash}p{0.27\linewidth}>{\raggedright\arraybackslash}X r@{}}
        \toprule
        Проверка & Проверяемые свойства & Результат \\
        \midrule
        Модульные тесты \texttt{ctest} &
        дерево, заголовок, битовый ввод-вывод, кодек, случайные данные,
        повреждённый вход & 6/6 \\
        Интеграционный сценарий \texttt{tests/test.sh} &
        полный цикл сжатие--распаковка и побайтовое сравнение файлов &
        6/6 \\
        \bottomrule
    \end{tabularx}
\end{table}

Интеграционные входы включают пустой файл, файл из повторяющегося символа,
текст, набор всех значений байта и случайные бинарные данные. Таким образом,
на вопрос RQ1 получен положительный ответ для предусмотренного набора
проверок: все запущенные тесты прошли успешно.

\subsection{Метрики и результаты измерений}

Сценарий \texttt{experiments/run\_all.sh} генерирует текстовые,
повторяющиеся, случайные и бинарно-подобные файлы размером 1, 10, 100~КиБ и
1~МиБ. Бинарно-подобный набор строится повторением системного исполняемого
файла, поэтому имитирует не случайный поток байтов, а фрагменты реального
бинарного файла. Кроме синтетических входов, сценарий добавляет типовые файлы:
объединённый файл исходного кода на C, PDF-файл отчёта и системный исполняемый
файл. Для каждого файла выполняется 10 запусков сжатия и распаковки. Основной
метрикой размера является отношение размера архива к размеру исходного файла;
значение меньше единицы означает уменьшение файла. Для времени приводятся
среднее значение и выборочное среднеквадратичное отклонение. В
таблице~\ref{tab:benchmark} приведены результаты для представительных входов.

\begin{table}[htbp]
    \centering
    \caption{Обработка представительных входов, 10 запусков; время указано в
    миллисекундах (меньше~--- лучше).}
    \label{tab:benchmark}
    \small
    \begin{tabular}{@{}l
            S[table-format=1.4]
            S[table-format=3.1(2)]
            S[table-format=3.1(2)]@{}}
        \toprule
        Вход & {\makecell{Отношение\\размеров}} &
        {\makecell{Сжатие,\\мс}} & {\makecell{Распаковка,\\мс}} \\
        \midrule
        Повторяющийся символ, 1~МиБ & 0.1270 & 115.0 \pm 2.7 & 83.7 \pm 1.6 \\
        Текст, 1~МиБ & 0.5064 & 138.6 \pm 1.4 & 107.6 \pm 2.3 \\
        Случайные байты, 1~МиБ & 1.0020 & 150.8 \pm 2.6 & 135.0 \pm 1.3 \\
        Исходный код C, 16~КиБ & 0.7744 & 64.7 \pm 1.0 & 64.8 \pm 1.0 \\
        PDF-отчёт, 65~КиБ & 1.0302 & 69.1 \pm 0.7 & 68.6 \pm 1.1 \\
        Системный бинарный файл, 151~КиБ & 0.4205 & 74.3 \pm 1.2 & 69.8 \pm 0.8 \\
        \bottomrule
    \end{tabular}
\end{table}

По вопросу RQ2 эксперимент показывает ожидаемую зависимость: файл с одним
доминирующим символом уменьшается приблизительно до 12,7\% исходного размера,
текстовый файл~--- до 50,6\%, а равномерно случайные байты практически не
сжимаются. Для малых файлов накладные расходы заголовка становятся
определяющими: архив случайного входа размером 1~КиБ имеет отношение размеров
2,9922. На типовых файлах результат различается: исходный код C сжимается до
77,4\% исходного размера, PDF-файл немного увеличивается из-за собственного
сжатия внутри формата, а системный бинарный файл уменьшается до 42,1\%. По
вопросу RQ3 для файлов размером 1~МиБ время сжатия в измеренном наборе
составляет от $115{,}0 \pm 2{,}7$ до $150{,}8 \pm 2{,}6$~мс, время
распаковки~--- от $83{,}7 \pm 1{,}6$ до $135{,}0 \pm 1{,}3$~мс.

\subsection{Ограничения воспроизводимости}

Сценарии эксперимента позволяют повторить измерения, но результаты не являются
бит-в-бит детерминированными. Файлы из набора \texttt{random} создаются из
\texttt{/dev/urandom}, поэтому при каждом запуске имеют другое содержимое.
Наборы \texttt{binary} и \texttt{real/system\_binary} строятся на основе
системного исполняемого файла \texttt{/bin/ls} или значения переменной окружения
\texttt{BINARY\_SOURCE}; следовательно, их содержимое зависит от операционной
системы и выбранного исходного файла, хотя размер синтетического бинарного
входа фиксируется скриптом.

Временные измерения зависят от процессора, нагрузки системы, версии компилятора
и режима сборки. Поэтому приведённые значения следует рассматривать как
характеристику конкретного окружения запуска, а не как переносимую оценку
производительности на любой платформе. Для повторения полного сценария также
требуются утилиты \texttt{python3} и \texttt{bc}. Пакеты Python
\texttt{pandas} и \texttt{matplotlib} нужны только для построения графиков: при
их отсутствии сценарий выводит таблицу агрегатов и пропускает генерацию
изображений. Дополнительным развитием эксперимента является
проверка на более широком наборе пользовательских файлов, например видеофайлов
и документов разных форматов.
